\section{Experiments}
\label{experiments}
% \vspace{-0.5em}

Based on performance and accessibility considerations, we choose the most recently published open-source video diffusion model, VideoCrafter~\citep{chen2024videocrafter2}, as our pre-trained video model in this paper. All experiments are conducted based on this model.

\noindent\textbf{Evaluation Metrics.} To evaluate video quality, we report Fréchet Video Distance (FVD) \citep{unterthiner2018towards}, Kernel Video Distance (KVD) \citep{kvd}. Since the tuning-free methods are supposed to keep the quality of the original pre-trained inference, we calculate the FVD and KVD between the original generated videos and videos generated by trajectory control methods. We use CLIP Similarity (CLIP-SIM) \citep{radford2021learning} to measure the semantic similarity among frames. In addition, we utilize the off-the-shelf detection model, OWL-ViT-large \citep{minderer2022simple}, to obtain the bounding box of the synthesized objects. Then Mean Intersection of Union (mIoU) and Centroid Distance (CD) are calculated to evaluate the trajectory alignment. CD is the distance between the centroid of the generated object and the input mask, normalized to $1$. When OWL-ViT-large fails to detect the target object in the generated videos, the farthest point will be assigned as the penalty in CD.

\subsection{Evaluation of Trajectory Control}

\begin{table}[t]
\caption{\textbf{Quantitative comparison of trajectory control.} FreeTraj achieves competitive performance in metrics about video quality and gains the best scores in metrics that are related to trajectory control.}
\vspace{-0.5em}
\label{tbl:score}
\begin{center}
\scalebox{0.8}{\begin{tabular}{lccccc}
\hline
\noalign{\smallskip}
\multicolumn{1}{c}{\bf Method}  &\multicolumn{1}{c}{\bf FVD ($\downarrow$)} &\multicolumn{1}{c}{\bf KVD ($\downarrow$)} &\multicolumn{1}{c}{\bf CLIP-SIM ($\uparrow$)} &\multicolumn{1}{c}{\bf mIoU ($\uparrow$)} &\multicolumn{1}{c}{\bf CD ($\downarrow$)}\\
\noalign{\smallskip}
\hline
\noalign{\smallskip}
Direct         &\textbf{118.19}  &\textbf{-2.28} & \textbf{0.980} & 0.161  & 0.225  \\
Peekaboo~\citep{jain2023peekaboo}       &403.00   &25.30   & 0.963  & 0.235  & 0.179  \\
TrailBlazer~\citep{ma2023trailblazer}    &556.00   &42.14  & 0.958  & 0.179  & 0.219  \\
Ours           &436.22   &29.85  & 0.956  &\textbf{0.281}  &\textbf{0.154}  \\
\noalign{\smallskip}
\hline
\end{tabular}}\vspace{-0.5em}
\end{center}
\end{table}

\begin{figure*}
\centering
\includegraphics[width=0.95\linewidth]{Figures/fig_trajcompare.pdf}
\vspace{-1.0em}
\caption{\textbf{Qualitative comparison of trajectory control.} We compare our proposed FreeTraj with direct inference (Direct), Peekaboo (Peek), and TrailBlazer (TraB). FreeTraj successfully generates high-fidelity results and is more accurate for trajectory control.}
\vspace{-1.0em}
\label{fig:trajcompare}
\end{figure*}


We compare our proposed FreeTraj to other tuning-free trajectory-controllable video generation methods with diffusion models Peekaboo~\citep{jain2023peekaboo} and TrailBlazer~\citep{ma2023trailblazer}. 

As shown in Figure~\ref{fig:trajcompare}, TrailBlazer has the worst control because it does not apply the control in spatial self-attention while the other two methods do. Videos generated by Peekaboo will roughly follow the given trajectory but not precisely. In addition, Peekaboo generates an additional black swan with weird artifacts, which is probably caused by the hard attention mask used in self-attention layers. Our FreeTraj succeeds in driving the target object following the given trajectories with vivid motions. 

For quantitative results shown in Table~\ref{tbl:score}, TrailBlazer has the worst mIoU and CD, which are only slightly better than those in direct inference. Peekaboo has better mIoU and CD, showing the rough control ability of trajectory but is still significantly weakened than our proposed FreeTraj. Both qualitative and quantitative results show that our FreeTraj is more accurate for trajectory control. However, we find the FVD, KVD, and CLIP-SIM, which are the references for video quality, are slightly worse than those in Peekaboo. As shown in Figure~\ref{fig:trajcompare}, Peekaboo tends to generate videos whose objects act around the center of the frame. This behavior is similar to reference videos which are directly generated by VideoCrafter, leading to a better FVD and KVD. TrailBlazer also mentions this phenomenon of lazy movement for better FVD.


In addition, we conducted a user study to evaluate our results based on human subjective perception. Participants were asked to watch the generated videos from all methods, with each example displayed in a random order to avoid bias. They were then instructed to select the best video in three evaluation aspects: trajectory alignment, video-text alignment, and video quality. The results, as shown in Table~\ref{tbl:user}, demonstrate that our approach outperforms the baseline methods by a significant margin, achieving the highest scores in all aspects. Notably, our method received nearly $70\%$ votes in terms of trajectory alignment. This user study confirms the superiority of our approach in terms of trajectory alignment, video-text alignment, and video quality.

\begin{table}[t]
\caption{\textbf{User study.} Users are requested to pick the best one among our proposed FreeTraj with the other baseline methods in terms of trajectory alignment, video-text alignment, and video quality.}
\vspace{-0.5em}
% 
\label{tbl:user}
\begin{center}
\scalebox{0.8}{\begin{tabular}{lccc}
\hline
\noalign{\smallskip}
\multicolumn{1}{c}{\bf Method}  &\multicolumn{1}{c}{\bf Trajectory Alignment} &\multicolumn{1}{c}{\bf Video-Text Alignment}
&\multicolumn{1}{c}{\bf Video Quality} \\
\noalign{\smallskip}
\hline
\noalign{\smallskip}
Peekaboo~\citep{jain2023peekaboo}       &6.48\%  &15.12\%  &13.58\% \\
TrailBlazer~\citep{ma2023trailblazer}   &8.03\%  &6.79\%  &6.79\% \\
Ours w/o Noise            &19.75\%  &15.43\%  &15.74\% \\
Ours                   &\textbf{65.74\%}  &\textbf{62.65\%}  &\textbf{63.89\%} \\
\noalign{\smallskip}
\hline
\end{tabular}}\vspace{-0.5em}
\end{center}
\end{table}

\subsection{Ablation Studies}

\begin{figure*}
\centering
\includegraphics[width=0.95\linewidth]{Figures/fig_trajabl.pdf}
\vspace{-1.0em}
\caption{\textbf{Ablation results.} (a) No noise guidance, (b) no high-frequency noise resampling, (c) hard attention mask, and (d) our whole method.}
\vspace{-1.0em}
\label{fig:trajabl}
\end{figure*}

\noindent\textbf{Ablation of Noise Guidance.}
%
To show the effectiveness of noise guidance, we run our designed attention guidance solely. Figure~\ref{fig:trajabl} (a) shows that pure attention guidance can also control the trajectory but may lose some accuracy. 

\noindent\textbf{Ablation of Attention Isolation.}
%
We also study two settings that may cause attention isolation. The first one is using the hard attention mask in Equation~\ref{eq:sata}. The second one uses no high-frequency noise resampling when applying trajectory injection in initial noises (Equation~\ref{eq:sata}).
%
Usually, diffusion models have some robustness to deal with the input with small deviation and recover it to generate qualified results. However, both of these two strategies will easily cause the value of the attention layer to deviate far from the data distribution in the training stage. It will lead to attention isolation where isolated regions almost pay no attention to other regions, losing the chance to recover back to the normal distribution. As shown in Figure~\ref{fig:trajabl} (b) and (c), blocky artifacts appear and follow the given trajectory in the generated videos. In addition, those artifacts happen to fall at the position of the attention mask or inject local noise.
